\documentclass[a4paper,11pt]{article}

\usepackage[utf8]{inputenc}
\usepackage[margin=0.8in]{geometry}
\usepackage{amsmath}
\usepackage{amssymb}
\usepackage{graphicx}
\usepackage{booktabs}
\usepackage{array}
\usepackage{tabularx}
\usepackage{caption}
\usepackage[hidelinks]{hyperref}
\usepackage{cite}
\usepackage{enumitem}
\usepackage{titlesec}
\usepackage{xcolor}
\usepackage{parskip}
\usepackage{float}

% House style: accent-coloured section headings with a rule underneath.
\definecolor{accent}{HTML}{0A66C2}
\titleformat{\section}{\large\bfseries\color{accent}}{\thesection}{0.6em}{}[\titlerule]
\titlespacing{\section}{0pt}{10pt}{6pt}
\titleformat{\subsection}{\normalsize\bfseries\color{accent}}{\thesubsection}{0.5em}{}
\titlespacing{\subsection}{0pt}{8pt}{4pt}

\title{\textbf{An Explainable Multi-Modal AI Framework for Maritime Security
Intelligence from Sentinel-1 SAR: Dark-Vessel Detection, Illegal-Fishing
Flagging, and Oil-Spill Assessment}}

\author{
  Shaun Marvell Rodrigues\\
  \textit{Supervisor:} Navneet Bhaskar\\[2pt]
  Nitte Mahalinga Adyantaya Memorial Institute of Technology (NMAMIT)
}
\date{}

\begin{document}

\maketitle
\thispagestyle{empty}

% -----------------------------------------------------------------------
\section*{Abstract}
% -----------------------------------------------------------------------

Maritime regions face escalating threats from illegal, unreported and unregulated
(IUU) fishing, ``dark'' vessels that disable their Automatic Identification System
(AIS) to evade monitoring, and deliberate or accidental oil discharge, yet these
problems are typically addressed by separate, stovepiped surveillance systems. In
this paper we present a single, end-to-end maritime-security framework built
entirely on free Sentinel-1 Synthetic Aperture Radar (SAR), which images through
cloud and darkness and is therefore suited to operational, all-weather
surveillance. From one georeferenced scene the system detects ships, segments oil
slicks, cross-checks detections against AIS to flag non-broadcasting vessels, and
rolls the outputs into interpretable security and environmental risk scores. Ship
detection couples a training-free, resolution-agnostic censored-mean Constant False
Alarm Rate (CFAR) detector with a YOLO11m oriented-bounding-box model trained on
HRSID (mAP@50 $=0.938$; $0.910$ after scale-augmented fine-tuning for robustness at
Sentinel-1's 10\,m/px resolution); the CFAR stage recovers vessels that the
high-resolution-trained detector misses at medium resolution, and a censoring step
corrects the classical CFAR target-masking failure in dense anchorages and near
ocean fronts. Oil-spill segmentation uses a SegFormer transformer benchmarked
against DeepLabV3+ and U-Net; on a strictly held-out test partition the
oil-class IoU is $\approx0.48$. We further document a train/inference parity lesson---a
whole-scene-resize inference path that resolves a $4\times$ resolution skew which had
silently driven oil predictions to zero---and the polarization asymmetry between
oil (cross-pol sensitive) and ships (polarization robust). Explainability is
provided through Grad-CAM saliency and rule-contribution bars rather than tabular
attribution. We discuss the cross-sensor and cross-region domain gaps that bound
generalization and an honest AIS workaround for the coverage limits of free
historical AIS. The result is an integrated, reproducible system that turns raw SAR
into actionable maritime-security intelligence; we note throughout that SAR yields
oil \emph{area/extent}, never volume.

\medskip
\noindent\textbf{Keywords:} Maritime Surveillance, Synthetic Aperture Radar,
Sentinel-1, Dark Vessel Detection, CFAR, Oil Spill Segmentation, SegFormer,
AIS Fusion, Explainable AI, Grad-CAM

% -----------------------------------------------------------------------
\section{Introduction}
% -----------------------------------------------------------------------

The world's oceans cover more than seventy percent of the planet and carry the
overwhelming majority of global trade, yet they remain among the least observed
environments on Earth. This observational gap is exploited at scale. Illegal,
unreported and unregulated (IUU) fishing is estimated to account for more than a
fifth of the global catch, depleting stocks and undermining the economies of
coastal states. Vessels engaged in IUU fishing, sanctions evasion, smuggling, or
unauthorised operations in restricted waters routinely disable or spoof their
Automatic Identification System (AIS) transponders to avoid being tracked, becoming
so-called \emph{dark} vessels that are present in the water but absent from the
cooperative reporting picture. At the same time, deliberate bilge-water discharge
and accidental spills release oil into the marine environment, threatening
fisheries, protected habitats and coastal communities. Each of these threats is
serious in isolation; together they describe a maritime-security problem that is
fundamentally about detecting activity that actively tries not to be seen.

Conventional maritime monitoring is poorly matched to this problem for three
reasons. First, it leans heavily on cooperative data: AIS is invaluable when present
but is trivially defeated by switching off the transponder, and recent ecological
work shows that AIS alone misses a large fraction of the vessels that satellite
imagery reveals inside marine protected areas~\cite{raynor2025}. Second, the optical
satellite imagery that the public associates with ``satellite surveillance'' is
blind at night and through cloud, and is typically months to years stale---unusable
for time-critical, all-weather operations. Third, the capabilities that do exist
tend to be deployed as separate, stovepiped systems---one tool for vessel detection,
another for fishing analysis, another for pollution monitoring---so that the
connections between them, such as linking a fresh oil slick to a nearby
non-broadcasting vessel, are never made automatically. Spaceborne Synthetic Aperture
Radar (SAR) addresses the first two limitations directly: as an active microwave
sensor it images day and night and penetrates cloud, and recent large-scale studies
have demonstrated that SAR can map industrial activity at sea that is invisible to
AIS~\cite{paolo2024}. What has been missing is an integrated, reproducible system
that turns a single SAR scene into a complete, explainable maritime-security
picture.

In this work we present such a system, built entirely on Sentinel-1, the
freely available European C-band SAR mission. Sentinel-1 is well suited to
operational use: its data are open, its Interferometric Wide-swath Ground Range
Detected (IW GRD) product images a 250\,km swath at roughly 10\,m/px on a six-day
revisit, and the constellation is being renewed by Sentinel-1C and Sentinel-1D as
the original satellites retire, so the C-band data format on which our pipeline
depends will persist. Because oil slicks and ships are extracted from the
\emph{same} georeferenced scene, they are automatically co-registered and
time-aligned, and AIS need only be matched to the scene's known acquisition
timestamp rather than to a live feed. From one scene the framework (i) detects every
vessel using a training-free censored-mean Constant False Alarm Rate (CFAR) detector
paired with a YOLO11m oriented-bounding-box network trained on the HRSID
dataset~\cite{wei2020}; (ii) segments oil slicks with a SegFormer~\cite{xie2021}
transformer trained on the publicly available Sentinel-1 SAR oil-spill dataset of
Trujillo-Acatitla~\textit{et~al.}~\cite{trujillo2024} (released on Zenodo), benchmarked
against DeepLabV3+ and U-Net; (iii) matches detections against historical AIS from Global Fishing Watch
to flag dark vessels; (iv) applies a rule-based zone-violation and risk-scoring layer
that links slicks to candidate responsible vessels; and (v) explains its outputs with
Grad-CAM saliency maps and rule-contribution bars. The complete pipeline is
orchestrated behind a service that emits georeferenced GeoJSON layers consumed by a
lightweight web dashboard for visualisation.

The contributions of this paper are deliberately framed around honest, reproducible
findings rather than headline benchmark numbers. (1)~We show that a censored-mean
CFAR detector provides \emph{resolution-agnostic} ship detection at Sentinel-1's
10\,m/px scale, where a detector trained on the much higher-resolution HRSID imagery
suffers a domain gap; the censoring step further corrects the classical CFAR
target-masking failure that suppresses dim vessels next to bright ships or ocean
fronts in dense anchorages. YOLO11m-OBB, used to confirm detections and supply
oriented geometry, reaches an mAP@50 of $0.938$ on HRSID and $0.910$ after
scale-augmented fine-tuning for medium-resolution robustness. (2)~We document a
train/inference parity lesson for oil segmentation: a whole-scene-resize inference
path that reproduces the training field of view and resolves a $4\times$ resolution
skew which had silently driven oil predictions to zero despite a healthy validation
score. (3)~We report a properly disjoint held-out evaluation in which the oil-class
IoU is $\approx0.48$, well below the in-distribution validation figure, and argue
that this validation-to-test gap---driven mainly by look-alike false positives---is
itself the scientifically useful result. (4)~We analyse the cross-sensor
(X/C/L-band) and cross-region domain gaps that bound generalisation and identify
xView3-SAR as the principled remedy for the ship-detection resolution gap. (5)~We
present an honest synthetic-AIS workaround for the coverage and latency limits of
free historical AIS. (6)~We deliver the whole pipeline as one operational system, not
a collection of disconnected models. Throughout, we are careful not to overclaim: a
recently published ``SOTA'' oil model whose code was never released falls back
silently to SegFormer and is reported as SegFormer; and we stress that SAR yields oil
\emph{area/extent}, never volume.

% -----------------------------------------------------------------------
\section{Related Work}
% -----------------------------------------------------------------------

Research relevant to this framework spans three threads: SAR vessel detection
(covering both classical Constant False Alarm Rate detectors and learned networks),
SAR oil-spill segmentation, and the fusion of SAR detections with AIS to reveal dark
vessels. We review ten representative works, emphasising publications from 2022 onward,
and position our design against each.

Paolo~\textit{et~al.}~\cite{paolo2022xview3} introduced xView3-SAR, the largest
labelled benchmark for SAR vessel and infrastructure detection, comprising nearly one
thousand analysis-ready Sentinel-1 scenes that are on average $29{,}400\times24{,}400$
pixels, together with bathymetry and wind-state layers. The associated international
challenge brought the problem of dark, AIS-absent vessels to the machine-learning
community and produced models that have since been deployed operationally against
illegal fishing. Crucially for our work, xView3 is Sentinel-1 Interferometric
Wide-swath imagery at roughly the same medium ($\sim$10--20\,m/px) resolution at which
our pipeline operates; we therefore identify fine-tuning on xView3-SAR as the
principled remedy for the resolution domain gap that affects our HRSID-trained
detector, and treat it as future work rather than a component of the present system.

Ma~\textit{et~al.}~\cite{ma2024amanet} proposed AMANet, an adaptive
multi-hierarchical attention network that fuses information across adjacent feature
layers and adaptively combines channel features, improving average precision on the
SSDD and HRSID benchmarks over its baseline. It is representative of the strong
attention-augmented CNN detectors that define recent HRSID performance. Like the great
majority of SAR ship-detection papers, however, AMANet is trained and evaluated at
HRSID's native $0.5$--$3$\,m/px resolution, where a $100$\,m vessel spans tens to
hundreds of pixels, and it does not address the distribution shift that arises at
$10$\,m/px Sentinel-1 IW resolution---the regime in which our pipeline must operate.

He~\textit{et~al.}~\cite{he2025acyolo} proposed AC-YOLO, a lightweight detector built
on YOLO11 that integrates an ACmix attention module, a CCFM feature-fusion neck, and an
MPDIoU regression loss, raising mAP by $1.5\%$ on HRSID and $1.2\%$ on SSDD while
reducing parameters by roughly a third. AC-YOLO confirms that YOLO11-class
architectures are the current high-water mark for SAR ship detection and directly
motivates our choice of a YOLO11m oriented-bounding-box backbone. As with AMANet,
however, the gains are demonstrated at training resolution; architectural refinement
alone does not close the medium-resolution inference gap, which is why our system pairs
the learned detector with a physics-based detector rather than relying on it alone.

Closest in spirit to our hybrid detector, Wen~\textit{et~al.}~\cite{wen2024cfaryolo}
embedded a sub-network that learns classical CFAR features inside YOLOv5s, enhancing
sensitivity to low-contrast targets. Our design differs in an important respect: rather
than folding CFAR into the learned detector as a trainable feature, we retain a
separate, \emph{training-free} censored-mean CFAR stage that is resolution-agnostic and
gates the YOLO detections, so the physics-based component requires no retraining when
the sensor or pixel size changes.

On the purely classical side, Zhang~\textit{et~al.}~\cite{zhang2023spcfar} developed an
improved superpixel-level CFAR (SP-CFAR) for high-resolution SAR, replacing the
rectangular clutter window with superpixels so that the detection threshold is
estimated from pure-clutter regions, which markedly improves inshore and densely packed
ship detection. SP-CFAR is effectively the non-deep-learning state of the art for dense
scenes, but a full superpixel pass over a multi-thousand-pixel scene is comparatively
slow and would conflict with the absolute-threshold land mask our pipeline relies on.
We therefore adopt censored-mean CFAR (CMLD), a faster member of the same robustness
family, which recovers most of SP-CFAR's benefit by censoring bright interferers from
the clutter estimate while preserving cell-averaging speed.

Turning to oil-spill segmentation, Wu~\textit{et~al.}~\cite{wu2024compositional}
proposed a compositional detector that pairs YOLOv8 with an adapted Segment Anything
Model, converting detected boxes into masks and fusing them through an ordered-mask
module. This SAM-based direction is precisely the lineage of the recently announced
``OilSAM2'' model we had hoped to benchmark; because that model's code was never
publicly released, our system instead uses a directly trained SegFormer, and we report
results honestly as SegFormer rather than as the unavailable state of the art.

Moon~\textit{et~al.}~\cite{moon2024diffusion} attacked the chronic data-scarcity
problem in SAR oil-spill segmentation using diffusion-generated synthetic imagery and
knowledge distillation over soft labels, improving minority-class performance without a
pre-trained teacher. Their approach is complementary to the look-alike hard-negative
upsampling we employ and points to a promising augmentation route for closing the
validation-to-test gap we report.

Most directly relevant to our central empirical finding,
Juarez~\textit{et~al.}~\cite{juarez2025crossdomain} studied cross-domain oil-spill
segmentation and reported that a model trained on Mediterranean data degrades from
$67.8\%$ to $51.8\%$ mIoU when transferred to a Peruvian domain, recovering up to
$+6$ mIoU with morphological region perturbation and synthetic SAR generation. Their
result independently corroborates the large cross-distribution drop we observe and
reinforces our argument that an honestly disjoint evaluation is essential when
reporting oil-segmentation accuracy.

On the fusion side, Paolo~\textit{et~al.}~\cite{paolo2024nature} mapped global
industrial activity at sea by combining Sentinel-1 SAR, AIS, and three deep
convolutional networks for object detection, length estimation, and vessel
classification, finding that $72$--$76\%$ of the world's industrial fishing vessels are
absent from public tracking. This study is the methodological anchor for the principle
that SAR reveals activity invisible to AIS; our framework operationalises the same
detect-then-cross-check logic at single-scene granularity and layers integrated risk
scoring and explainability on top.

Complementarily, Raynor~\textit{et~al.}~\cite{raynor2025science} combined AIS with SAR
to study fishing inside strictly protected marine areas and found that AIS missed close
to $90\%$ of the fishing vessels that SAR detected in coastal zones. This quantifies
precisely why AIS-only monitoring is insufficient and motivates the dark-vessel
flagging and zone-violation logic at the heart of our risk layer.

Taken together, these works advance individual components---learned or classical vessel
detection, oil-spill segmentation, and AIS--SAR fusion---but none integrates detection,
oil segmentation, AIS-based dark-vessel discrimination, rule-based risk scoring, and
visual explainability into a single operational pipeline driven by one freely available
sensor. That integration, together with the honest systems lessons we document, is the
contribution of this paper.

\end{document}
